\documentclass[10pt]{article}
\usepackage[margin=20mm]{geometry}
\usepackage{amsmath,amssymb,mathtools,fontspec,unicode-math,xurl,hyperref}
\IfFontExistsTF{DejaVu Serif}{\setmainfont{DejaVu Serif}}{\setmainfont{Latin Modern Roman}}
\setmathfont{latinmodern-math.otf}
\setlength{\parindent}{0pt}\setlength{\parskip}{6pt}\sloppy
\title{A scoped verification of curvature and log-quasimode growth}
\author{AI-assisted technical review}
\date{}
\begin{document}
\maketitle
Abstract: This report checks the interpolation and summation mechanism underlying the two logarithmic gains, the spectral-window sharpness argument, and selected boundary and exponent calculations in Huang and Sogge’s paper. The calculations reproduce the stated gains and the limsup formulation of sharpness. A direct flat-torus construction provides an independent special-case check. A local flow-box notation requires the signed intersection-time interpretation stated by the surrounding geometric construction; it does not furnish a counterexample to the main estimates. No substantive mathematical defect is established in the inspected scope. The global kernel and full bilinear estimates are dependencies, not independently reproved results of this review. Scientific adjudication remains pending.\par
\section*{Identity, version, and inspected scope}
The reviewed work is Xiaoqi Huang and Christopher D. Sogge, Curvature and sharp growth rates of log-quasimodes on compact manifolds, Inventiones mathematicae 239 (2025), 947–1008, DOI https://doi.org/10.1007/s00222-025-01315-2 [HS]. The acquired 62-page PDF is the publisher version of record. In this report, PDF page 1 is printed page 947, so a printed page is its PDF page plus 946.\par
Sections 1–4 were read from the acquired PDF’s extracted text. The independent calculations focus on Theorems 1.1–1.2; Proposition 2.4 and equations (2.50)–(2.71); the orthogonal spectral-window argument on printed page 973; and the use of Lemmas A.1–A.2 on printed pages 986–988. Selected original page images were checked to resolve mathematical notation. The detailed global parametrix estimates in Lemma 2.5, the full variable-coefficient bilinear argument in the Appendix, and the complete Bieberbach-group periodization are not independently established here. No separate supplement or released computation was required for these analytic checks.\par
\section*{The two logarithmic gains}
Let \(n\ge2\), \(q_c=2(n+1)/(n-1)\), and \[\mu(q)=\frac{n-1}{2}\left(\frac12-\frac1q\right),
 \qquad 2<q\le q_c,\qquad T=c_0\log\lambda.\] Theorem 1.1 assumes a connected compact Riemannian manifold with all sectional curvatures nonpositive, or strictly negative for its stronger conclusion. The exponent \(q\) is fixed; the negative-curvature constant may depend on \(q\).\par
Equation reference: eq:dyadic. The endpoint operator bounds used in Proposition 2.4 are \[\left\lVert W_N\right\rVert_{2\to2}\lesssim T^{-1}N,\qquad
 \left\lVert W_N\right\rVert_{1\to\infty}\lesssim
 T^{-1}\lambda^{(n-1)/2}N^{1-\alpha},\] on the indicated product of counting and Riemannian measure spaces, with dyadic \(2\le N\le T\) (printed pp. 969–970). Interpolating with weight \(\vartheta=1-2/q\) on the second bound gives \[
 \left\lVert W_N\right\rVert_{q'\to q}
 \lesssim T^{-1}\lambda^{2\mu(q)}
 N^{\,1-\alpha(1-2/q)}.\] Thus the sum is uniformly bounded in \(T\), after the factor \(T^{-1}\), when \(\alpha>q/(q-2)\). Taking the square root in the \(UU^*\) estimate yields \(\lambda^{\mu(q)}T^{-1/2}\), exactly (2.52). At \(q=q_c\), the exponent of \(N\) is \(1-2\alpha/(n+1)\). Summing the geometric series gives \[\left\lVert U\right\rVert_{2\to q_c}\lesssim\lambda^{1/q_c}
 \begin{cases}
 T^{-\alpha/(n+1)},&\alpha<(n+1)/2,\\
 T^{-1/2},&\alpha>(n+1)/2.
 \end{cases}\] The nonpositive-curvature value \(\alpha=(n-1)/2\) therefore gives \(T^{-1/q_c}\). The arbitrary fixed decay order in the strictly negative case allows \(\alpha>q/(q-2)\) for each fixed \(q>2\). These are precisely the different gains claimed.\par
The equality case \(\alpha=q/(q-2)\) would instead leave a dyadic factor \(\log T\) before the square root. Proposition 2.4 correctly excludes it. Likewise, \(q=2\) is not included, and letting \(q\downarrow2\) does not supply a uniform constant. The proof does not obtain the negative-curvature subcritical improvement merely by interpolating the critical estimate with the trivial \(L^2\) bound: it uses (2.45), then (2.52), as required.\par
For the Appendix’s high-ceiling calculation, put \(\bar q=2(n+2)/n\). Direct simplification gives \[(q_c-\bar q)\left(\frac{3(n-1)}{16}-\frac18\right)
 =\frac{3n-5}{4n(n-1)}>0,\quad n\ge2.\] For \(n=2\) it equals \(1/8\), consistent with the separate treatment on printed page 988. At \(q=\bar q\), (A.12) gives a power saving of \(3/(4n)-\varepsilon\) relative to \(\lambda^{\mu(\bar q)\bar q}\); hence it supplies the endpoint that is not stated in (A.13). These checks verify the displayed exponent arithmetic and absorption margins, conditional on the bilinear input estimates themselves.\par
\section*{Sharpness, quantifiers, and an independent flat model}
The definition immediately before Theorem 1.2 (printed p. 951) explicitly uses \[f=\Theta(g)\quad\Longleftrightarrow\quad
 0<\limsup_{\lambda\to\infty} f(\lambda)/g(\lambda)<\infty .\] It is not a lower bound at every sufficiently large frequency. Empty spectral windows, assigned value zero in the theorem, therefore do not contradict it.\par
The square-root window factor on printed page 973 follows from orthogonality. Partition a unit frequency interval into \(J\lesssim\delta^{-1}\) disjoint spectral pieces \(I_j\) of length at most \(\delta\). For \(K=\max_j\left\lVert\chi_{I_j}\right\rVert_{2\to q}\), \[\left\lVert\chi_{[\lambda,\lambda+1]}f\right\rVert_q
 \le K\sum_j\left\lVert\chi_{I_j}f\right\rVert_2
 \le K\sqrt J\,\left\lVert f\right\rVert_2 .\] Consequently sharp unit-band lower bounds imply a subsequential \(\delta^{1/2}\lambda^{\mu(q)}\) lower bound for some narrow windows. This is the precise implication used by the paper, taking the cited unit-band sharpness theorem as an external input. It does not prove a pointwise lower bound at every window.\par
An explicit model independently verifies the flat scaling. On \(\mathbb T^n=(\mathbb R/2\pi\mathbb Z)^n\) with normalized Haar measure, define \[F_{k,m}(x)=(2m+1)^{-(n-1)/2}e^{ikx_1}
       \prod_{j=2}^{n}\sum_{r=-m}^{m}e^{irx_j}.\] Its \(L^2\) norm is exactly one. Every frequency satisfies \[0\le\sqrt{k^2+|r'|^2}-k
 \le\frac{(n-1)m^2}{2k}.\] Choose \(m=\lfloor\sqrt{k\delta/(n-1)}\rfloor\) with \(k\delta\to\infty\). The spectrum then lies in \([k,k+\delta]\). On the transverse box \(|x_j|\le c/m\), with fixed sufficiently small \(c>0\), all summands have uniformly positive real parts. Hence \[\left\lVert F_{k,m}\right\rVert_q\gtrsim
 m^{(n-1)(1/2-1/q)}
 \asymp(k\delta)^{\mu(q)}.\] For \(n=2\), \(k=256\), \(\delta=1/4\), and \(m=8\), exact finite arithmetic gives \(\left\lVert F\right\rVert_{4}^{4}=193/17\), while the largest frequency increment is at most \(1/8\). This finite check is illustrative; the preceding analytic argument is the asymptotic justification. Neither calculation replaces the paper’s construction for arbitrary compact flat quotients.\par
\section*{A local notation clarification and its limits}
Equation (2.26), printed page 961, labels the time at which a geodesic meets \(\Pi=\{x_n=0\}\) by \(s=d_g(x,\Pi)\). In ordinary distance notation this cannot literally specify the hitting time for generic directions. In the Euclidean two-dimensional model, \(x=(0,h)\), \(h>0\), and a unit covector \(\xi=(\epsilon,\sqrt{1-\epsilon^2})\) with \(0<\epsilon<1\) give \((\Phi_h(x,\xi))_2=h+h\sqrt{1-\epsilon^2}>0\). The signed intersection time is instead \(s=-h/\xi_n\).\par
The preceding transverse-section construction and the explicit invariance requirement (2.27) identify the intended definition: extend the initial cutoff constantly along each geodesic, using its unique local intersection with \(\Pi\). Transversality \(\xi_n\ne0\) makes that hitting time smooth and yields (2.27). This replacement changes no theorem statement or decay exponent. This report records the local notation clarification and does not treat it as evidence that the main microlocal construction or the estimates fail. It also does not certify every subsequent symbol calculation merely from this correction.\par
\section*{Conclusion}
The checked interpolation, summation, spectral-window, and model sharpness calculations support the stated conclusions within their hypotheses. No substantive defect is established within this scope. In particular, the limsup convention, the strict inequalities in Proposition 2.4, and the fixed-\(q\) constants remove several apparent endpoint objections.\par
Dimension | Recorded state\par
Assessment | No Issue Found within the inspected scope; one local notation clarification\par
Repair verification | Signed flow-intersection interpretation checked locally; no general theorem repair claimed\par
Scientific decision | Pending\par
Human review / adoption | No human review supplied; author adoption not requested\par
Remaining limits | Global kernel estimates, full bilinear proof, general flat-quotient construction, and external sharpness theorem not independently reproved\par
\section*{Supporting checks}
The symbolic checks verify identities under the stated domains; they are not a proof of the paper’s analytic estimates. It is an initial scientific delivery, not an independent acceptance or a human endorsement. Exact execution identities and hashes are recorded in provenance rather than inferred from expected values.\par
1 X. Huang and C. D. Sogge, Curvature and sharp growth rates of log-quasimodes on compact manifolds, Invent. Math. 239 (2025), 947–1008. Publisher version of record, published online 29 January 2025. https://doi.org/10.1007/s00222-025-01315-2.\par
\end{document}
